% Source: https://artofproblemsolving.com/community/c3h1565700_factorization
% Stable identity: aops-topic-1565700-q01
\textbf{1.} The square root $\sqrt{2xy}$ is a real number exactly when $xy\ge0$. In that case, by the difference of squares,
\[
\left(x+y-\sqrt{2xy}\right)\left(x+y+\sqrt{2xy}\right)=(x+y)^2-2xy=x^2+y^2 .
\]
So the identity holds for all real $x,y$ with $xy\ge0$ (both nonnegative or both nonpositive), and it is meaningless over the reals when $xy<0$.

\textbf{2.} Suppose $x^2+y^2=P(x,y)\,Q(x,y)$ with $P$, $Q$ nonconstant real polynomials. Total degrees add, so both have degree $1$: $P=ax+by+e$, $Q=cx+dy+f$. Comparing the coefficients of $x^2$ gives $ac=1$, so $a\ne0$. Now put $y=1$:
\[
x^2+1=\bigl(ax+(b+e)\bigr)\bigl(cx+(d+f)\bigr),
\]
and the first factor vanishes at the real number $x=-\dfrac{b+e}{a}$. But $x^2+1>0$ for every real $x$, a contradiction. Hence $x^2+y^2$ has no factorisation into real polynomials.

The student's factors are not polynomials: $\sqrt{2xy}$ is not a polynomial in $x$ and $y$ (and it is not even defined for $xy<0$). ``$x^2+y^2$ does not factor'' is a statement about polynomial factors, so there is no contradiction. (Over the complex numbers it does factor: $x^2+y^2=(x+iy)(x-iy)$.) $\blacksquare$
